% =====================================================================
%  CONCLUSIONES (versión recortada, para comparar con
%  capitulos/09-conclusiones.tex)
%
%  Un párrafo por objetivo específico y uno por la pregunta problema, como
%  sugiere la plantilla. Revisión del 27 de octubre: cada párrafo dice qué se
%  concluye del objetivo, no repite los resultados; las cifras quedan en el
%  capítulo de resultados, salvo la cota de la tasa de éxito. Sin primera
%  persona y sin remisiones a otras secciones.
% =====================================================================
\chapter*{Conclusiones}
\addcontentsline{toc}{chapter}{CONCLUSIONES}

El primer objetivo específico, modelar la arquitectura funcional del sistema, se
cumplió. Un coordinador central con un agente por nivel convierte el cambio de piso en un problema de coordinación y no de locomoción. Ningún agente necesita
subir escaleras ni usar el ascensor, y lo que se traspasa en la escalera es la
responsabilidad de guiar al usuario. Separar la planificación del relevo de
ROS~2 permitió verificarla de forma exhaustiva y sin simulador, y esa
verificación encontró defectos que las corridas no habrían mostrado. La
asignación de tareas resulta trivial por la estructura del problema, con un solo
agente admisible por nivel; la dificultad está en coordinar dos agendas que
dependen una de otra durante el relevo, y ahí está el aporte del trabajo.

El segundo objetivo, la plataforma robótica con los dos agentes, se cumplió. El
AWS DeepRacer sirve como agente de un sistema de guiado en interiores, con tres
condiciones que este trabajo hizo explícitas. Su odometría depende del entorno:
sin codificadores en las ruedas, el avance solo se mide si el pasillo tiene
estructura, y por eso la elección del sitio de operación es parte del diseño y
no un detalle de la prueba. Su tracción no responde por debajo de 0,40~m/s, lo
que fija la precisión de llegada en 0,5~m. Y su tarjeta de dos núcleos obliga a
repartir el cómputo, con la navegación a bordo y la coordinación en un
computador aparte. Queda abierta la estabilidad en sesiones de varias horas.

El tercer objetivo, la interfaz de interacción con el usuario, se cumplió. Una
página web servida en la red local basta para pedir, seguir, confirmar y
cancelar una misión desde cualquier teléfono, sin instalación y sin salida a
internet. La confirmación de cambio de piso y la cancelación, añadidas al ver el
sistema en uso, son las funciones que más lo acercan a una operación real. La primera impide dar por terminada una misión con el usuario en el piso
equivocado, y la segunda devuelve al agente a su escalera en cualquier etapa.

El cuarto objetivo, la evaluación experimental, se cumplió en simulación y con
los agentes físicos. La campaña en simulación permite afirmar, con 95~\% de
confianza, que la tasa de éxito del sistema supera el 70~\%, y la campaña física
muestra que el sistema completo funciona en un edificio real con la misma cota.
Los fallos de las dos campañas estuvieron en la precisión de llegada de los
agentes y no en la coordinación. La evaluación dejó además una conclusión de
método: decidir qué mide cada criterio fue más difícil que construir el sistema.
Un criterio que no puede fallar, un reloj sin resolución para la magnitud que
mide o un modelo del edificio sin verificar producen resultados que parecen evidencia y no lo son. Ninguno de esos casos habría aparecido corriendo más experimentos.

Sobre la pregunta problema, cómo asistir la orientación de usuarios dentro de
entornos interiores con múltiples pisos utilizando sistemas robóticos, se
concluye que un esquema de relevo entre dos agentes, uno por nivel, es una
respuesta viable. Los agentes se reparten el guiado de una persona que cambia
de piso sin que ninguno atraviese la zona de transición y sin modificar el
edificio. Queda sin medir la experiencia de la persona guiada, que corresponde a
un estudio con usuarios.
